\documentclass[a4paper, 11pt]{article}
\usepackage[top=3cm, bottom=3cm, left = 2cm, right = 2cm]{geometry}
\geometry{a4paper}
\usepackage{ottalt}
\usepackage{amsmath,amssymb,amsthm}
\usepackage{mathpartir}
\theoremstyle{remark}
\newtheorem*{remark}{Remark}
\title{A formal description of the type system}
\author{notch1p}
\date{August 21, 2026}
\inputott{rules.ott}

\newenvironment{rules}{\bigskip}{\bigskip}

\begin{document}
\maketitle
\section{Introduction}

The purpose of this document is to provide readers with a up-to-date formal
description of Tigris' type system, with a heavy emphasis on its implementation
-- judgements are all algorithmic, specifically.
It is advised that this be read alongside all the notes\footnote{
  Readers are also assumed familiarity with the various works this document quotes, for
  being the sturdy foundation that Tigris relys on.
}
in
\texttt{typing/} to form a more complete understanding of ``how we did it'',
``why it is done this way''. Meanwhile, System F elaboration is trivial modulo instance
resolution and dictionary synthesis -- since the description is not about the latter, we omit the former
here completely.

For an \emph{much outdated} review and formal description of the language,
I suggest reading doc/midterm (in English) for a work-in-progress yet self-contained report,
or my B.Eng. thesis%
\footnote{TinyTheorem: A dependently-typed programming language with light theorem proving (I really should name it \emph{Tigris} in the title.)}
(in Chinese) located in the repo (and the submodule in \texttt{doc/}) \texttt{beng-thesis}.
Directories \texttt{examples/} and \texttt{tests/cases/} showcase some of the
(mostly) up-to-date features and capability of the language.

\paragraph{Layout.} We don't follow the traditional formal grammar - judgements order but
instead gives every important datatype their own sections.

Below we present the core calculus at both term and type level. Not everything is presented here,
for brevity. Refer to \texttt{syntax.ott} for the complete syntaxes.

\section{Terms}
\nonterms{expr}

let-expressions and toplevel bindings have been simplified to take single binding. Same
applies to case analysis. A recursive binding is no different from a normal one except
its \textbf{fun}-shaped RHS is wrapped in a \textbf{fix} node with its name as the head argument; Others are processed as-is.
Literals are omitted.

The Let(rec)-group in the surface language is inferred in SCC order, thus a nonrecursive
binding may still be mutually recursive iff each one of them does not mention itself. Note that
this is different from OCaml's Let-group semantics.

\section{Types}
\nonterms{tyvar,typeexpr,scheme,predicate}
Tigris is in System F. The type system supports arbitrary ranked types,
allows polytypes to nest everywhere in the types\footnote{Though a higher-rank scheme cannot carry predicates, for now.}.
Later we'll show how \emph{partial} impredicativity is obtained ``for free''
without Quick Look (Serrano et al. ICFP'20), unlike GHC -- and why much of what we
once boasted was unsound.

The usually-absent type lambda is part of the syntax directly to encode both higher-kinded
types and type synonyms (abbreviations). It is handled with care to ensure no beta-redexes
survive construction whereas partial applicated type constructors get eta-expanded. Unlike GHC,
we permit partially applicated type synonyms.

Typeclasses in the surface language is syntactic sugar (with additional infomation recorded) for structures,
which, in turn, is syntactic sugar for single-constructor inductive types whose type constructor's \emph{name}
is the same as the value constructor. This invariant is enforced by the parser. A predicate is a typeclass constraint.
Unlike GHC, typeclass constructor (often higher-rank) can be called directly and the methods
can be explicitly pattern matched out of a instance with generalization.

We now describe expression typing. Note that toplevel binding is omitted since in practice
we simply wrap them in a Letexp whose body is \textsf{Unit}.

\subsection{Constraint Generation.} Tigris implements a constraint solver with wanted pools
similar to that of OutsideIn(X) based on the understanding that type
inference is a process of generating and solving constraints (Pottier, \'{R}emy).

\begin{remark}
  We deem a judgement \emph{most} ``algorithmic'' if it generates constraints. Those
  judgements uses $\vDash$ whereas others use $\vdash$ unless told otherwise. Readers
  should not assume judgements are ``pure'' since most of them carry a state for
  e.g. counters used for fresh type variables.
\end{remark}

\nonterms{constraint}

\paragraph{Term Inference.} $[[R ; G |= e => ty | Cs]]$ reads,
under typing environment $[[G]]$ and rigid set $[[R]]$,
$e$ is inferred to be $ty$. Below we give the trivial ones directly, then elaborate
on mroe interesting cases.

\drules[I]{$[[R ; G |= e => ty | Cs]]$}{term inference}{Var,Lam,Prod,Ascribe-Mono,Cond,Fix,Match}

\subparagraph{Function Application.} We eagerly consume the arrow when possible. In the
case that the function's instantiated type is not (yet) an arrow, possibly
a metavariable, we leave an ordinary equation to the pool. A forall that surfaces in it only
after solving can only unify with another forall -- it is never instantiated (see Instantiation). A rank-n
function hiding behind a metavariable thus needs an ascription guiding it, e.g. with
$g : \textsf{Int}\to(\forall\alpha.\,\alpha\to\alpha)\to\textsf{Int}$, \texttt{apply g 1 id}
needs $(\forall\alpha.\,\alpha\to\alpha)\to\textsf{Int}$ on \texttt{apply g 1}.
The two cases are given below as

\begin{mathpar}
  \drule{I-App-Arrow}
  \drule{I-App-MV}
\end{mathpar}

Type ascriptions are divided into rank-1 and rank-n. The shallow subsumption $\vdash_\textsf{sh}$
from jones2007 gave rise to two derivation \rref{Spec,Skol}
which corresponding to the two \texttt{solveAll}s in rank-1. Meanwhile, we switch to
check mode for rank-n modulo predicates for which the paper's \rref{Gen2,Annot} provide
guidances.
Monotype ascription is trivial enough that we've given it above already.

\subparagraph{Rank-1 Ascription.} \Rref{Spec} solves against the instantiated
annotation whereas \rref{Skol} solves against its skolemization, with escape-check
on rigids. It is worth noting the solving here is \emph{lax}, or \emph{best effort},
in the comments, which simply means ``solve what we can'': by
returning a partial subst, leaving the rest to the outer pool
(i.e. the outer call to \texttt{solveAll}) to try again instead of signaling an error.
Readers may see this as a kind of backtracking -- we try again in hopes that the outer pool
may reveal more solutions that makes the unsolvables solvable.

\begin{mathpar}
  \drule{I-Ascribe-RankOne}
  \drule{I-Ascribe-RankN}
\end{mathpar}

\subparagraph{Rank-N Ascription.} An interesting case is that we do not have prenex
conversion\footnote{
  the coercion counterpart of deep skolemization $\vdash_\textsf{dsh}$.
  The latter is used for \emph{annotated} lambdas \rref{AAbs2}
  which is absent here at syntax level. Though we need them back if we were to support OCaml's
  locally abstract types or GHC's scoped type variables. Not a real concern at the moment.
}
that \rref{Gen2} requires. From what we have gathered, a direct result seems to be, that, e.g.
the two supposedly isomorphic types are no longer so,
\[
  \forall \alpha\,\beta.\,\alpha\to\beta\to\beta \ncong \forall \alpha.\,
  \alpha\to\forall \beta.\,\beta\to\beta.
\]
This used to be papered over by \rref{U-Inst1}, which turned out to be unsound and has been
removed. Now ascribing the latter to a variable $f$ of the former type is
rejected though eta-expansion ($\lambda x.\,f\,x$ is accepted) circumvents it -- which aligns
with GHC's behavior\footnote{The author verified this claim on his machine with GHC 9.14.1.}.

\subparagraph{Let-expression.}
``Let should not be generalized'' is a polemic about abolishing local let generalization for type system
with complex extensions. While our system isn't complex enough to justify that, the author's exposition
about generalizing in OutsideIn(X) remains relevant to us as well: A to-be-generalized free var shared by both
pending predicates and result types would be erroneously monomorphized by solve-time
resolution the pool introduces, missing what could be generalized.

\begin{mathpar}
  \drule{I-Let}
\end{mathpar}
where $L$ denotes the constraint set that contains \emph{leftover} equations;
\textsf{rigidize} and \textsf{unskolem} kept opaque for brevity.

OutsideIn(X) proposes the concept of ``untouchables'' to bar refinement from touching those variables.
While we could do some plumbing in the pool and maintain a so-called \emph{untouchable set}, the
straightfoward, readily available tool in our system is \emph{skolemization}. Thus, typing a Let-expression
requires the following steps:

\begin{enumerate}
  \item Solves the group's equations to exhaustion with refinement disabled. Which always makes sense
    whether we take predicate constraint into consideration or not.
  \item Skolemizes the abstract-able\footnote{
      The exact condition is given in the rules; A crucial point is that
      we must not quantify over free vars in $L$ since they are determined
      by the pool, not \emph{robbed} by us;
      Also, ``abstract'' and ``generalize'' are used interchangably.
    }
    free vars on top of the substituted result from step 1.
  \item Re-solve again with refinement enabled.
  \item Manually consume the group's constraint to avoid
    outer pools re-solving the same, unskolemized constraints, otherwise it defeats
    the very work we've done. Given the fact that the pool is non-consuming in the sense it does not
    consumes constraints by itself. Practically, \texttt{solveAll} isn't even in the typechecking
    monad \texttt{InferC}, pure-by-construction.
  \item Finally, undo step 2 (unskolemize) and generalize by plain HM(X).
\end{enumerate}

\paragraph{Checking.}
Checking monotypes is trivial. It is more interesting to see what unrestricted
binding under checking allows us to achieve.

\drules[C]{$[[R ; G |= e <= ty | Cs]]$}{term checking}{Skol,Default,Lam-Mono,Lam-Poly}
\Rref{C-Skol} is also a derivation of the shallow subsumption.
This way we can make sure the type of $[[e]]$
is as general (if not more) as that of the user.

\Rref{Lam-Mono,Lam-Poly} are the keys to arbitrary-rank types. Concretely in \rref{Lam-Poly}
the parameter is bound with the domain as-is, a scheme, so that a polymorphic parameter
instantiates (\rref{I-Var}) at each use in the body.

Impredicativity, on the other hand, comes with the representation. Since we allow a metavariable
to bind anything (polytypes, specifically), a type variable may be instantiated with a polytype
-- the impredicativity, by definition, of System F. For example, consider
\begin{align*}
  \texttt{run} &: (\forall s.\,s\to s)\to\textsf{Int},\\
  \texttt{id}\,\textsf{run} &: ((\forall s.\,s\to s)\to\textsf{Int})\to((\forall s.\,s\to s)\to\textsf{Int}),\\
  \texttt{Cons} &: (\forall\alpha.\,\alpha\to\alpha)\to\textsf{List}\,(\forall\alpha.\,\alpha\to\alpha)\to\textsf{List}\,(\forall\alpha.\,\alpha\to\alpha),
\end{align*}
where we instantiated quantified TVs with polytypes.
Note the difference: a polytype at the top is aliased to a type variable only from a type, as in the latter,
never from a term, since \rref{I-Var,C-Default} instantiate it first -- so much for \textit{free}.
However, such a solution is reached by equations alone
(\rref{U-Var-Bind,U-TSch}), and a polytype that is only known after solving -- hidden behind a
metavariable when the term was inferred or checked -- stays opaque. It is neither instantiated nor
skolemized, and becomes usable only once it is syntactically visible again: through let-generalization,
by substituting it in, or through an ascription; Plain program would get rejected\footnote{
  see \texttt{examples/impred} for concrete programs.
}.
Polytypes in constructor fields, e.g.
$\texttt{PCons} : (\forall\alpha.\,\alpha\to\alpha)\to\textsf{PList}\to\textsf{PList}$, are
unaffected, since the field's type is visible both when constructing and when matching.

Notice the claim is quantified with \textit{partial}. The deleted \rref{U-Inst1} instantiated any forall against
non-polys in an equation thus a polytype behind a metavariable \emph{was} instantiated at its uses:
\texttt{(id run) id}, \texttt{apply g 1 id}, $[\texttt{id}] : \textsf{List}\,(\forall\alpha.\,\alpha\to\alpha)$, and matching an element out of such a list to
apply it all typechecked. That part of impredicativity was obtained erroneously which we shall
explain below.

\paragraph{Instantiation.}
A central part of our typechecker, which, is especially the case after first-order refactoring,
unification is no longer higher-order -- homogenous unification for polytypes up to renaming
by \rref{U-TSch}. Instantiation thus happens in
\rref{I-Var,I-App-Arrow,C-Default} and \rref{Spec} of rank-1 ascription, where we know
the provided and expected types.

\begin{remark}[Why \rref{U-Inst1} is unsound]
  There used to be one more instantiation, \rref{U-Inst1}, inside unification: a scheme
  $[[forall </ tvi // i /> . ty]]$ against a monotype $[[ty']]$ instantiated $[[</ tvi // i />]]$
  with fresh metavariables and unified the body with $[[ty]]$.
  \begin{mathpar}
    \drule{U-InstOne}
  \end{mathpar}
  It was meant for rank-n functions
  hiding behind a metavariable e.g. \texttt{apply g 1 id}, on the grounds that skolemizing instead
  would reject $\forall\alpha.\,\alpha\to\alpha \sim \beta\to\textsf{Bool}$, which solves with
  $\alpha\mapsto\textsf{Bool}$. That argument -- $\vdash\mathsf{i}\sigma\Rightarrow\tau$ holds only for
  \begin{enumerate}
    \item a poly on the \emph{provided} side,
    \item the check is at the top of that provided/expected comparison.
  \end{enumerate}
  but binrel $(\sim)$ is commutative -- equation commutes and \rref{U-Inst1}
  isn't guaranteed to be at the top -- we satisfies neither.
  Specifically, unification is symmetric (though not shown, the rules are implicitly commutative),
  and \rref{U-Arr,U-Prod,U-App} recurse into arrow domains first (where the direction flips)
  and type arguments before the rule is reached. To put it simply, \rref{U-Inst1} instantiated polytypes that \textit{should've stayed poly}
  -- the solution (substs) need not equate both sides
  -- the one thing unification promises.
  If we were to keep it, we would unsoundly accept:
  \begin{itemize}
    \item $h : (\textsf{Int}\to\textsf{Int})\to\textsf{Int}$ aliased to
      $k : (\forall\alpha.\,\alpha\to\alpha)\to\textsf{Int}$, rank-2 only: $k$ applies its argument
      to \texttt{true}, and $h$ passes it a \texttt{(1 + \_)};
    \item \texttt{apply g 1} $(\lambda x.\,x+1)$, the very case the rule was for, with a $g$
      that applies its argument to \texttt{true};
    \item \texttt{x ++ y} where $x : \textsf{List}\,(\textsf{Int}\to\textsf{Int}),\, y : \textsf{List}\,(\forall\alpha.\,\alpha\to\alpha)$: the pool first solves the metavariable
      of \texttt{(++)} to the polytype, then the rule accepts
      $\textsf{List}\,(\textsf{Int}\to\textsf{Int}) \sim \textsf{List}\,(\forall\alpha.\,\alpha\to\alpha)$
      under $\textsf{List}$.
  \end{itemize}
  It is now removed at the expense of impredicativity discussed under Checking\footnote{
    \texttt{error/\{rankn-\{contra,mv\},impred-share\}} have more counterexamples.
  }.
\end{remark}

\end{document}